\subsection*{The opposite-prime Fox target for the third-power cover}

\paragraph{TF.1 (Proved): a marked full-chain target with odd deck group and characteristic-two coefficients.}
Let $F_5=C'$ with the fixed five cube-Schreier generators
$(g_1,\ldots,g_5)$, and set
\[
 D_3=F_5'F_5^3=C''(C')^3,\qquad
 A_3=F_5/D_3\simeq(C_3)^5,\qquad
 k=\mathbb F_4=\mathbb F_2[\omega],\quad
 \omega^2+\omega+1=0.
 \tag{TF.1a}
\]
Thus $[F_5:D_3]=243$, $[C:D_3]=1944$, and $D_3$ is
free profinite of rank $1+243(5-1)=973$. The five-petal
covering graph has the exact sequence
\[
 0\longrightarrow H_1(D_3,k)\longrightarrow k[A_3]^5
 \stackrel{\partial}{\longrightarrow}I_{A_3}
 \longrightarrow0,\qquad
 \partial=(a_1-1,\ldots,a_5-1),\quad a_j=\bar g_j.
 \tag{TF.1b}
\]
Its dimensions are $(973,1215,242)$. The finite-index
subgroup and graph may first be constructed in the discrete
free group; completion preserves the finite-coefficient
abelianization of this free subgroup.

Use the same based path convention as FF.1:
\[
 d(uv)=d(u)+\bar u\,d(v),\qquad d(g_j)=e_j,
                     \qquad\partial d(w)=\bar w-1.
 \tag{TF.1c}
\]
The actual finite path group is constructed over
$\mathbb F_2$, before scalar extension to $k$:
\[
 F_5/(D_3'D_3^2)\simeq
 \{(v,a)\in\mathbb F_2[A_3]^5\rtimes A_3:
                              \partial v=a-1\}.
 \tag{TF.1d}
\]
This follows because a path to each deck sheet exists,
and arbitrary graph cycles supply all possible additional
chains. On $D_3$, the kernel is exactly its mod-two
abelianization kernel $D_3'D_3^2$. Extension to $k$ is
injective and does not enlarge the actual source group.

Only the deck index $243$ is averaged in the Fourier
construction, and it is a unit in characteristic two.
For each character $\chi:A_3\to\langle\omega\rangle$,
put
\[
 \partial_\chi=(\chi(a_1)-1,\ldots,\chi(a_5)-1),
 \qquad V_\chi=k^5,\qquad M_\chi=\ker\partial_\chi.
\]
The trivial character has $M_1=V_1$ of dimension five.
For each of the $242$ nontrivial characters,
$\partial_\chi\ne0$, $M_\chi$ has dimension four, and
\[
             0\longrightarrow k^4\longrightarrow k^5
                     \stackrel{\partial_\chi}{\longrightarrow}k
                     \longrightarrow0
 \tag{TF.1e}
\]
is retained in full. The homology dimensions are therefore
$973=5+242\cdot4$.

For an actual source automorphism $\phi$ induced by $L$,
let $J_\phi$ have columns $d(\phi(g_j))$. The chain rule
and marked transport are unchanged:
\[
 \begin{aligned}
 T_\phi(v)&=J_\phi\bar\phi(v),
 &\qquad \eta&=\chi\circ\bar\phi^{-1},\\
 T_\phi|_{V_\chi}&=J_\phi(\eta):V_\chi\longrightarrow V_\eta,
 &\partial_\eta J_\phi(\eta)&=\partial_\chi.
 \end{aligned}
 \tag{TF.1f}
\]
In particular all words are multiplied with their actual
character permutations. The full sum of the nontrivial
blocks is an $L$-stable finite representation, denoted
$\rho_3$. These statements follow from the literal Fox
chain rule, not an assumed five-point GT action.
No division by $[C:D_3]=1944$ or by $[C:F_5]=8$ occurs.

\paragraph{TF.2 (Proved): exact faithfulness on the required derived-source restriction.}
For every $w\in F_5'=C''$,
\[
 \boxed{\quad
 \rho_3(w)=1\quad\Longleftrightarrow\quad
                  w\in D_3'D_3^2.\quad}
 \tag{TF.2a}
\]
Equivalently,
\[
       \ker(\rho_3|_{F_5'})=F_5'\cap(D_3'D_3^2).
 \tag{TF.2b}
\]
The intersection is part of the statement: $D_3'D_3^2$
is not being asserted to lie in $F_5'$.
There is no assertion here of faithfulness on all of $C$,
or of centrelessness of its characteristic-two quotient.

Indeed $w\in F_5'$ first implies $w\in D_3$, so its deck
coordinate is one. The exact inner Fox formula gives
\[
       T_{w,\chi}=I-d(w)_\chi\partial_\chi
                =I+d(w)_\chi\partial_\chi
                         \quad(\operatorname{char}k=2).
 \tag{TF.2c}
\]
For every nontrivial character this matrix is the identity
if and only if $d(w)_\chi=0$, because $\partial_\chi$
has a nonzero entry. The remaining trivial component is
already zero: $d(w)_1$ is the vector of exponent sums
modulo two, and $w\in F_5'$. Fourier inversion therefore
makes all these identities equivalent to $d(w)=0$.
By (TF.1d), this is exactly $w\in D_3'D_3^2$.
This proves both directions of (TF.2a).

In particular the retained RA.5 pentagon defect belongs
to $D_2'\subset F_5'$, so it satisfies this hypothesis.
For that actual class of witnesses,
\[
 \boxed{\quad
 \rho_3(\mathcal P(f))=1
 \quad\Longleftrightarrow\quad
 \mathcal P(f)=1\text{ in }L/(D_3'D_3^2).\quad}
 \tag{TF.2d}
\]
The quotient is genuine because $D_3'D_3^2$ is
characteristic in $D_3$, itself characteristic in the
normal subgroup $C'$ of $L$.
The final defect is blockwise because it belongs to
$F_5'$; its individual pentagon factors need not be.

Writing $m_\chi=d(w)_\chi$, its residual blocks have the
precise form
\[
 \partial_\chi m_\chi=0,\qquad
 \operatorname{rank}(T_{w,\chi}-I)\leq1,
 \qquad (T_{w,\chi}-I)^2=0,
 \qquad T_{w,\chi}^2=I.
 \tag{TF.2e}
\]
For any $j$ with $\chi(a_j)\ne1$, the vector is recovered
from the full chain matrix by
\[
                  m_\chi=
             \frac{e_j-T_{w,\chi}e_j}{\chi(a_j)-1}.
\]
The $121$ conjugate pairs satisfy
$m_{\chi^2}=m_\chi^2$ coefficientwise. Hence the new
translation data have dimension $242\cdot4=968$ over
$\mathbb F_2$ after descent from $\mathbb F_4$.
In fact the image of $F_5'$ in $F_5/(D_3'D_3^2)$ is
exactly this elementary abelian group of rank $968$.
To see this, the finite extension by the odd-order
abelian deck group $A_3$ splits by the same factor-set
averaging proof as FF.5. Its derived subgroup is therefore
$[H_1(D_3,\mathbb F_2),A_3]$, the sum of its nontrivial
character components. This is a statement about all
words in $F_5'$, not a claim that all its vectors occur
as pentagon defects of genuine gentle elements.

\paragraph{TF.3 (Proved): the same restricted comparison at every two-adic level.}
Let $\mathcal O_2=\mathbb Z_2[\omega]$ be the unramified
quadratic extension with $\omega^2+\omega+1=0$, and
use all nontrivial full Fox blocks over
$\mathcal O_2/2^n$. For every $n\geq1$ and $w\in F_5'$,
\[
 \boxed{\quad
        \rho_{3,n}(w)=1
        \quad\Longleftrightarrow\quad
                    w\in D_3'D_3^{2^n}.\quad}
 \tag{TF.3a}
\]
Indeed $243$ remains a unit, the differences of distinct
cube roots are units, and the original path group is
constructed over $\mathbb Z/2^n$ before faithful scalar
extension. The nontrivial boundary rows therefore still
have unit entries. The exponent-sum component of $w$
is exactly zero, so the proof of TF.2 applies verbatim.
This detects the complete two-primary abelian component
on the stated restriction. It does not assert that the
other prime components are tested, that the quotient of
$C$ is centreless, or that a source automorphism extends
merely from one pentagon identity. Those are separate
conditions if an action or an adelic conclusion is sought.
